\subsection{Fundamentos técnicos}

\subsubsection{Conceptos arquitectónicos: nodos y \textit{topics}}

Un sistema basado en ROS se organiza como un grafo de procesos que se
comunican entre sí, denominado grafo de ROS. La unidad básica de este grafo es
el \textbf{nodo}: un participante que realiza una función lógica específica,
como leer un sensor, estimar la posición del robot o controlar un motor. Cada
nodo utiliza una biblioteca cliente para comunicarse con el resto del sistema:
\texttt{roscpp} y \texttt{rospy} en ROS~1, y \texttt{rclcpp} y \texttt{rclpy}
en ROS~2, para los lenguajes C++ y Python respectivamente. Los nodos pueden
comunicarse con otros nodos alojados en el mismo proceso, en otro proceso o en
otro equipo de la red. Las conexiones entre ellos se establecen mediante un
mecanismo de descubrimiento, cuya implementación difiere entre ROS~1 y ROS~2,
según se detalla más adelante.

El mecanismo de comunicación fundamental es el \textbf{\textit{topic}}: un canal
con nombre sobre el que opera un modelo de publicación/suscripción. Los nodos
productores de datos (\textit{publishers}) publican mensajes en un
\textit{topic}, y todos los nodos suscritos a él (\textit{subscribers}) los
reciben. Un mismo \textit{topic} admite cero o más publicadores y cero o más
suscriptores, lo que configura un esquema de comunicación de muchos a muchos.
Este modelo presenta dos propiedades relevantes:
\begin{itemize}
    \item \textbf{Anonimato:} el suscriptor no necesita conocer qué nodo
    originó el dato. Esto permite reemplazar, agregar o quitar nodos sin
    modificar el resto del sistema.
    \item \textbf{Tipado fuerte:} cada \textit{topic} transporta un único tipo
    de mensaje, definido en un archivo \texttt{.msg}. A partir de esta
    definición, las herramientas de ROS generan automáticamente el código para
    los distintos lenguajes soportados.
\end{itemize}

Los \textit{topics} resultan adecuados para flujos de datos continuos, como
las lecturas de sensores o el estado del robot, donde no se espera una
respuesta por parte del receptor.

\subsubsection{Patrones de comunicación: \textit{services} y \textit{actions}}

Además de los \textit{topics}, ROS ofrece dos patrones de comunicación
orientados a la interacción entre un cliente y un servidor.

Un \textbf{\textit{service}} implementa una llamada a procedimiento remoto
(RPC, \textit{Remote Procedure Call}) bajo un esquema de solicitud y
respuesta. Un nodo cliente envía una solicitud, el nodo servidor realiza el
procesamiento correspondiente y devuelve un resultado. Su interfaz se define en
un archivo \texttt{.srv}, compuesto por un mensaje de solicitud y uno de
respuesta. Cada nombre de \textit{service} debe tener un único servidor, aunque
puede ser utilizado por cualquier cantidad de clientes. Dado que el cliente
generalmente queda a la espera del resultado, los \textit{services} están
pensados para operaciones de ejecución breve, como consultar un estado o
modificar una configuración.

Una \textbf{\textit{action}} es una llamada a procedimiento remoto de larga
duración, que incorpora la posibilidad de recibir retroalimentación durante la
ejecución y de cancelar el objetivo antes de su finalización. Un ejemplo típico
es la orden de desplazar el robot hasta una posición determinada: la tarea
puede demorar segundos o minutos, durante los cuales el servidor informa el
progreso y el cliente puede cancelarla. Su interfaz se define en un archivo
\texttt{.action}, compuesto por tres mensajes: objetivo (\textit{goal}),
resultado (\textit{result}) y retroalimentación (\textit{feedback}).

En ROS~2, las \textit{actions} no constituyen un mecanismo de transporte
independiente, sino que se construyen sobre los dos patrones anteriores: cada
\textit{action} se compone internamente de tres \textit{services} (envío del
objetivo, cancelación y obtención del resultado) y dos \textit{topics}
(retroalimentación y estado). En ROS~1, en cambio, las \textit{actions} no
formaban parte del núcleo y se implementaban mediante una biblioteca
independiente denominada \texttt{actionlib}.

\subsection{Ámbitos de aplicación}

Donde se utiliza y como, en que mercados, etc.

\subsection{Estrategias Competitivas}
¿Competencia? ¿Producto único?

\subsection{Normativas asociadas}\label{sec:normativas}
¿Qué normas tiene asociadas ITU, IEEE?